\documentclass[10pt,a4paper]{amsart}
\usepackage[margin=19mm]{geometry}
\usepackage{amsmath,amssymb,mathtools}
\usepackage[scr=rsfs,cal=euler]{mathalfa}
\usepackage{xfrac}
\usepackage[unicode,hidelinks,bookmarks=false]{hyperref}
\newcommand{\ie}{i.e.\ }
\newcommand{\cf}{cf.\ }
\newcommand{\resp}{respectively\ }
\newcommand{\Subsection}{\subsection}
\providecommand{\nobreakdash}{\nobreak\hbox{-}}
\setlength{\emergencystretch}{2em}
\multlinegap0pt
\usepackage{amssymb}
\usepackage{dsfont}
\usepackage{mathtools}
\usepackage[shortlabels]{enumitem}
\newtheorem{thm}{Theorem}
\title{Effective correlation and decorrelation for newforms, and weak subconvexity for $L$-functions}
\author{Nawapan Wattanawanichkul}
\date{Source: arXiv:2405.05249v2 (CC BY 4.0)}
\begin{document}
\maketitle

\noindent\emph{Source and licence.} Effective correlation and decorrelation for newforms, and weak subconvexity for $L$-functions. Version arXiv:2405.05249v2 (CC BY 4.0), distributed by arXiv under the Creative Commons Attribution 4.0 International licence, http://creativecommons.org/licenses/by/4.0/, as recorded in the arXiv licence metadata for this exact version and shown on the abstract page https://arxiv.org/abs/2405.05249v2. Full licence notice: \texttt{licenses/arxiv-2405.05249-CC-BY-4.0.txt}.\par\smallskip

\noindent\textit{Editorial scope.} Counted unit: the article's thm thm:1 (source0.tex lines 1747--1753). The article's complete original proof of it is delivered with it (source0.tex lines 1754--1800, one \texttt{proof} environment of 47 lines). No mathematics is written, repaired or re-derived: the statement and the proof are the article's own lines, in the article's own notation, and no retained line is substituted or re-flowed.  No further source-local context is needed: every cross-reference of every retained line resolves inside the two delivered spans.  Nothing was omitted from any retained span, so the package carries no omission record.  No retained line contains a citation, so the wrapper carries no bibliography and no bibliography entry.  No retained line contains a figure environment, an image-inclusion command or drawing code, so no image is delivered and the package declares no asset dependency.  Editorial formatting only, in addition: an AMS \texttt{amsart} preamble that restates the article's own declarations and macros used by the retained lines, namely the environments \texttt{thm} on separate counters, together with the standard packages those lines need.  The retained environments print in this document as Theorem 1 in order of appearance, whereas the article prints the same items with the numbering of its own full document.  The complete native source of the article as distributed in the arXiv e-print for this exact version is bundled unchanged as \texttt{sources/158865/source0.tex}.
\begin{thm}
\label{thm:1}
Let $f$ be a holomorphic cuspidal Hecke eigenform and $\phi$ be a Hecke--Maa{\ss} newform, each of level $1$.  There exists an absolute constant $0<\omega\leq \frac{1}{10}$ and a function $H(s)$ such that in the region $\mathrm{Re}(s)\geq \frac{1}{2}-\omega$, $H(s)$ is holomorphic, $H(s)\asymp 1$, and
\[
L(s,\mathrm{Ad}\, f\times \phi)=H(s) L(2s,\mathrm{Ad}\, f) L(2s,\mathrm{Ad}\, \phi) \sum_{p\mid n\implies p\ge 19}\frac{\lambda_{f}(n^2)\lambda_{\phi}(n)}{n^{s}}.
\]
\end{thm}

\begin{proof}
By multiplicativity of the Hecke eigenvalues, it suffices to work locally.  Let $p$ be prime.  For any $\pi\in\mathfrak{F}_2$ of conductor $1$, we have that
\begin{align*}
L_p(s,\pi) &= \prod_{\ell\in\{-1,1\}}\frac{1}{1-\alpha_{\pi}(p)^{\ell} p^{-s}}=\sum_{j=0}^{\infty}\frac{\lambda_{\pi}(p^j)}{p^{js}},\\
L_p(s,\mathrm{Ad}\, \pi) &= \prod_{\ell\in\{-1,0,1\}}\frac{1}{1-\alpha_{\pi}(p)^{2\ell} p^{-s}}=\sum_{j=0}^{\infty}\frac{\lambda_{\mathrm{Ad}\, \pi}(p^j)}{p^{js}}.
\end{align*}
We write $\alpha_p = \alpha_{f}(p)$, $\beta_p = \alpha_{\phi}(p)$, and
\[
L_p(s,\mathrm{Ad}\, f \times \phi) = \prod_{\ell \in\{-1,0,1\}}\,\prod_{\ell'\in\{-1,1\}}\frac{1}{1-\alpha_p^{2\ell} \beta_p^{\ell'} p^{-s}}=\sum_{j=0}^{\infty}\frac{\lambda_{\mathrm{Ad}\, f\times\phi}(p^j)}{p^{js}}.
\]
Note that $\lambda_f(p^j)= \sum_{r=0}^j \alpha_p^{2r-j}$ and $\lambda_{\phi}(p^j)= \sum_{r=0}^j \beta_p^{2r-j}$.  By \eqref{eqn:GRC2} and the Ramanujan conjecture for $\pi$, we have the bounds
\begin{equation}
\label{eqn:GRC2}
|\alpha_p|=1,\qquad \max\{|\beta_p|,|\beta_p|^{-1}\}\leq p^{\frac{7}{64}}.
\end{equation}

Consider the rational function
\begin{equation}
\label{eqn:aux_rational}
\frac{1+(\beta_p+\beta_p^{-1})p^{-s}-(\alpha_p^2+1+\alpha_p^{-2})p^{-2s}}{(1-\alpha_p^2 \beta_p p^{-s})(1-\alpha_p^2 \beta_p^{-1}p^{-s})(1-\alpha_p^{-2} \beta_p p^{-s})(1-\alpha_p^{-2} \beta_p^{-1}p^{-s})}.
\end{equation}
On one hand, we can expand \eqref{eqn:aux_rational} as a polynomial in $p^{-s}$ times a product of four geometric sums.  On the other hand, we can identify \eqref{eqn:aux_rational} as a product of local $L$-functions.  Identifying these two expressions, we conclude that if $\mathrm{Re}(s)>7/64$, then
\begin{equation}
\label{eqn:key}
L_p(s,\mathrm{Ad}\,f\times\phi)=\frac{L_p(s,\phi)}{1+\lambda_{\phi}(p)p^{-s}-\lambda_{f}(p^2)p^{-2s}}\sum_{j=0}^{\infty}\frac{\lambda_{f}(p^{2j})\lambda_{\phi}(p^j)}{p^{js}},
\end{equation}
provided that the denominator for the leading factor on the right-hand side is nonzero.  We verify that if $p\geq 19$ and $\mathrm{Re}(s)\geq \frac{2}{5}$, then $|1+\lambda_{\phi}(p)p^{-s}-\lambda_{f}(p^2)p^{-2s}|>0.06$ using \eqref{eqn:GRC2} and Mathematica 12.  For such $p$ and $s$, we find that
\begin{equation}
\label{eqn:G_def}
\frac{L_p(s,\phi)}{1+\lambda_{\phi}(p)p^{-s}-\lambda_{f}(p^2)p^{-2s}}=1+(\lambda_{f}(p^2)+\lambda_{\phi}(p^2))p^{-2s} +\cdots,
\end{equation}
which agrees with the beginning of the expansion $L_p(2s,\mathrm{Ad}\,f)L_{p}(2s,\mathrm{Ad}\,\phi)$ up to an error of $O(p^{-3s})$. This leads us to define
\[
H_p(s) = \frac{L_p(s,\phi)}{(1+\lambda_{\phi}(p)p^{-s}-\lambda_{f}(p^2)p^{-2s})L_p(2s,\mathrm{Ad}\, f)L_p(2s,\mathrm{Ad}\,\phi)}.
\]
By our above numerical calculation and \eqref{eqn:GRC2}, we know that if $p\geq 19$,  then in the region $\mathrm{Re}(s)\geq \frac{2}{5}$, $H_p(s)$ is holomorphic and
\begin{align*}
L_p(s,\mathrm{Ad}\,f\times\phi) &= H_p(s)L_p(2s,\mathrm{Ad}\, f)L_p(2s,\mathrm{Ad}\,\phi)\sum_{j=0}^{\infty}\frac{\lambda_{f}(p^{2j})\lambda_{\phi}(p^j)}{p^{js}},\\
H_p(s)&=1-\lambda_{f}(p)^2\lambda_{\phi}(p)p^{-3s}+(\lambda_{f}(p)^2(\lambda_{\phi}(p)^2+1)-2)p^{-4s}+\cdots.
\end{align*}

Since the coefficients of $p^{-s}$ and $p^{-2s}$ in the expansion of $H_p(s)$ are zero, \eqref{eqn:GRC2} implies that there exists an absolute constant $0<\omega\leq\frac{1}{10}$ such that if $\mathrm{Re}(s)\geq \frac{1}{2}-\omega$, then the functions
\[
\mathcal{H}_1(s)=\prod_{p\geq 19}H_p(s),\qquad \mathcal{H}_2(s) = \prod_{p<19}\frac{L_p(s,\mathrm{Ad}\,f\times\phi)}{L_p(2s,\mathrm{Ad}\,f)L_p(2s,\mathrm{Ad}\,\phi)}
\]
are holomorphic and $\asymp 1$.  To finish, we choose $H(s) = \mathcal{H}_1(s)\mathcal{H}_2(s)$.
\end{proof}

\end{document}
